\documentclass[10pt]{extarticle}
\usepackage[a4paper,margin=1.5cm]{geometry}
\usepackage[T1]{fontenc}
\usepackage{amsmath,amssymb}
\usepackage{graphicx}
\usepackage[export]{adjustbox}
\usepackage{multicol}
\usepackage{enumitem}
\usepackage{booktabs,array,tabularx}
\setlength{\columnsep}{0.7cm}
\setlength{\parindent}{0pt}
\setlength{\parskip}{3pt}
\pagestyle{empty}
\makeatletter
\newcommand{\qfigdef}[4]{\@namedef{qf@file@#1}{#2}\@namedef{qf@w@#1}{#3}\@namedef{qf@wide@#1}{#4}}
\newcommand{\qf@check}[1]{\@ifundefined{qf@file@#1}{\PackageError{score2}{Unknown figure id #1}{}}{}}
\newcommand{\qf@img}[4][]{\includegraphics[width=#2\dimexpr\@nameuse{qf@w@#4}pt\relax,max width=#3,max height=0.45\textheight,#1]{\@nameuse{qf@file@#4}}}
\newif\ifqf@full
\newif\ifqf@boxed
\newsavebox{\qf@box}
\newcommand{\qf@open}{\global\qf@fulltrue\par\global\qf@boxedtrue
  \setbox\qf@box\vbox\bgroup\hsize\textwidth\linewidth\textwidth}
\newcommand{\qf@emit}[1]{\par\egroup\global\qf@boxedfalse
  \ifdim\dimexpr\ht\qf@box+\dp\qf@box\relax<#1\textheight\noindent\usebox{\qf@box}\par\else\unvbox\qf@box\fi}
\newcommand{\qfullsplit}{\ifqf@boxed\qf@emit{0.85}\fi}
\newcommand{\qf@close}{\ifqf@boxed\qf@emit{0.45}\fi\global\qf@fullfalse}
\newcommand{\qfullbegin}{\end{multicols}\qf@open}
\newcommand{\qfullend}{\qf@close\begin{multicols}{2}}
\newcommand{\qfullstart}{\qf@open}
\newcommand{\qfulldone}{\qf@close}
\newcommand{\qfig}[2][1]{\qf@check{#2}%
  \ifqf@full\qf@blk{#1}{#2}\else
  \ifnum\@nameuse{qf@wide@#2}=1 \expandafter\@firstoftwo\else\expandafter\@secondoftwo\fi
  {\end{multicols}\begin{center}\qf@img{#1}{\textwidth}{#2}\end{center}\begin{multicols}{2}}%
  {\qf@blk{#1}{#2}}\fi}
\newcommand{\qf@blk}[2]{\par\smallskip{\centering\qf@img{#1}{\linewidth}{#2}\par}\smallskip}
\newcommand{\qfigi}[2][1]{\qf@check{#2}\qf@img[valign=c]{#1}{\linewidth}{#2}}
\makeatother
\newcommand{\papersec}[1]{\par\bigskip{\centering\Large\bfseries PAPER-#1\par}\medskip}
\newcommand{\partsec}[2]{\par\medskip\noindent{\large\bfseries PART-#1: #2}\par\smallskip}
\newcommand{\sect}[1]{\par\smallskip\noindent\textbf{#1}\par}
\newcommand{\qn}[1]{\par\smallskip\noindent\textbf{#1.}~}
\newcommand{\paraq}[2]{\par\smallskip\noindent\textbf{Paragraph for Questions #1 and #2}\par}
\newenvironment{qoptions}{\begin{itemize}[leftmargin=1.8em,itemsep=1pt,topsep=2pt,label={}]}{\end{itemize}}
\newcommand{\opt}[2]{\item[(#1)] #2}
% figure registry (generated by score2_tools.py finalize)
\qfigdef{p1_phy_s1_q02_1}{JEE Advanced (2021 Series) - Test 05 (11 Apr 2021)-img/p1_phy_s1_q02_1.png}{305.5}{1}
\qfigdef{p1_phy_s1_q03_1}{JEE Advanced (2021 Series) - Test 05 (11 Apr 2021)-img/p1_phy_s1_q03_1.png}{55.5}{0}
\qfigdef{p1_phy_s1_q04_1}{JEE Advanced (2021 Series) - Test 05 (11 Apr 2021)-img/p1_phy_s1_q04_1.png}{168.5}{0}
\qfigdef{p1_phy_s2_q08_1}{JEE Advanced (2021 Series) - Test 05 (11 Apr 2021)-img/p1_phy_s2_q08_1.png}{171.5}{0}
\qfigdef{p1_phy_s2_q09_1}{JEE Advanced (2021 Series) - Test 05 (11 Apr 2021)-img/p1_phy_s2_q09_1.png}{120.5}{0}
\qfigdef{p1_phy_s3_q01_1}{JEE Advanced (2021 Series) - Test 05 (11 Apr 2021)-img/p1_phy_s3_q01_1.png}{165.4}{0}
\qfigdef{p1_phy_s3_q02_1}{JEE Advanced (2021 Series) - Test 05 (11 Apr 2021)-img/p1_phy_s3_q02_1.png}{109.3}{0}
\qfigdef{p1_phy_s3_q03_1}{JEE Advanced (2021 Series) - Test 05 (11 Apr 2021)-img/p1_phy_s3_q03_1.png}{77.0}{0}
\qfigdef{p1_phy_s3_q05_1}{JEE Advanced (2021 Series) - Test 05 (11 Apr 2021)-img/p1_phy_s3_q05_1.png}{87.5}{0}
\qfigdef{p1_phy_s3_q06_1}{JEE Advanced (2021 Series) - Test 05 (11 Apr 2021)-img/p1_phy_s3_q06_1.png}{142.3}{0}
\qfigdef{p1_chem_s1_q02_1}{JEE Advanced (2021 Series) - Test 05 (11 Apr 2021)-img/p1_chem_s1_q02_1.png}{284.3}{1}
\qfigdef{p1_chem_s1_q02_2}{JEE Advanced (2021 Series) - Test 05 (11 Apr 2021)-img/p1_chem_s1_q02_2.png}{164.1}{0}
\qfigdef{p1_chem_s1_q02_3}{JEE Advanced (2021 Series) - Test 05 (11 Apr 2021)-img/p1_chem_s1_q02_3.png}{164.1}{0}
\qfigdef{p1_chem_s1_q02_C}{JEE Advanced (2021 Series) - Test 05 (11 Apr 2021)-img/p1_chem_s1_q02_C.png}{153.8}{0}
\qfigdef{p1_chem_s1_q02_D}{JEE Advanced (2021 Series) - Test 05 (11 Apr 2021)-img/p1_chem_s1_q02_D.png}{153.8}{0}
\qfigdef{p1_chem_s1_q04_1}{JEE Advanced (2021 Series) - Test 05 (11 Apr 2021)-img/p1_chem_s1_q04_1.png}{238.0}{0}
\qfigdef{p1_chem_s1_q05_1}{JEE Advanced (2021 Series) - Test 05 (11 Apr 2021)-img/p1_chem_s1_q05_1.png}{84.9}{0}
\qfigdef{p1_chem_s1_q05_2}{JEE Advanced (2021 Series) - Test 05 (11 Apr 2021)-img/p1_chem_s1_q05_2.png}{69.6}{0}
\qfigdef{p1_chem_s1_q05_3}{JEE Advanced (2021 Series) - Test 05 (11 Apr 2021)-img/p1_chem_s1_q05_3.png}{69.6}{0}
\qfigdef{p1_chem_s2_q08_B}{JEE Advanced (2021 Series) - Test 05 (11 Apr 2021)-img/p1_chem_s2_q08_B.png}{32.9}{0}
\qfigdef{p1_chem_s2_q08_C}{JEE Advanced (2021 Series) - Test 05 (11 Apr 2021)-img/p1_chem_s2_q08_C.png}{56.3}{0}
\qfigdef{p1_chem_s2_q08_D}{JEE Advanced (2021 Series) - Test 05 (11 Apr 2021)-img/p1_chem_s2_q08_D.png}{69.8}{0}
\qfigdef{p1_chem_s2_q10_1}{JEE Advanced (2021 Series) - Test 05 (11 Apr 2021)-img/p1_chem_s2_q10_1.png}{110.8}{0}
\qfigdef{p1_chem_s2_q11_A}{JEE Advanced (2021 Series) - Test 05 (11 Apr 2021)-img/p1_chem_s2_q11_A.png}{30.4}{0}
\qfigdef{p1_chem_s2_q11_B}{JEE Advanced (2021 Series) - Test 05 (11 Apr 2021)-img/p1_chem_s2_q11_B.png}{45.4}{0}
\qfigdef{p1_chem_s2_q11_C}{JEE Advanced (2021 Series) - Test 05 (11 Apr 2021)-img/p1_chem_s2_q11_C.png}{71.4}{0}
\qfigdef{p1_chem_s2_q11_D}{JEE Advanced (2021 Series) - Test 05 (11 Apr 2021)-img/p1_chem_s2_q11_D.png}{44.7}{0}
\qfigdef{p1_chem_s3_q04_1}{JEE Advanced (2021 Series) - Test 05 (11 Apr 2021)-img/p1_chem_s3_q04_1.png}{43.5}{0}
\qfigdef{p1_chem_s3_q05_1}{JEE Advanced (2021 Series) - Test 05 (11 Apr 2021)-img/p1_chem_s3_q05_1.png}{226.9}{0}
\qfigdef{p1_chem_s3_q05_3}{JEE Advanced (2021 Series) - Test 05 (11 Apr 2021)-img/p1_chem_s3_q05_3.png}{245.2}{0}
\qfigdef{p1_chem_s3_q05_5}{JEE Advanced (2021 Series) - Test 05 (11 Apr 2021)-img/p1_chem_s3_q05_5.png}{230.0}{0}
\qfigdef{p1_chem_s3_q05_7}{JEE Advanced (2021 Series) - Test 05 (11 Apr 2021)-img/p1_chem_s3_q05_7.png}{161.3}{0}
\qfigdef{p1_chem_s3_q05_9}{JEE Advanced (2021 Series) - Test 05 (11 Apr 2021)-img/p1_chem_s3_q05_9.png}{158.0}{0}
\qfigdef{p2_phy_s1_q01_1}{JEE Advanced (2021 Series) - Test 05 (11 Apr 2021)-img/p2_phy_s1_q01_1.png}{182.8}{0}
\qfigdef{p2_phy_s1_q03_A}{JEE Advanced (2021 Series) - Test 05 (11 Apr 2021)-img/p2_phy_s1_q03_A.png}{132.0}{0}
\qfigdef{p2_phy_s1_q03_B}{JEE Advanced (2021 Series) - Test 05 (11 Apr 2021)-img/p2_phy_s1_q03_B.png}{132.5}{0}
\qfigdef{p2_phy_s1_q03_C}{JEE Advanced (2021 Series) - Test 05 (11 Apr 2021)-img/p2_phy_s1_q03_C.png}{131.8}{0}
\qfigdef{p2_phy_s1_q03_D}{JEE Advanced (2021 Series) - Test 05 (11 Apr 2021)-img/p2_phy_s1_q03_D.png}{133.5}{0}
\qfigdef{p2_phy_s1_q04_1}{JEE Advanced (2021 Series) - Test 05 (11 Apr 2021)-img/p2_phy_s1_q04_1.png}{127.8}{0}
\qfigdef{p2_phy_s1_q05_1}{JEE Advanced (2021 Series) - Test 05 (11 Apr 2021)-img/p2_phy_s1_q05_1.png}{194.4}{0}
\qfigdef{p2_phy_s1_q06_1}{JEE Advanced (2021 Series) - Test 05 (11 Apr 2021)-img/p2_phy_s1_q06_1.png}{153.4}{0}
\qfigdef{p2_phy_s2_q03_1}{JEE Advanced (2021 Series) - Test 05 (11 Apr 2021)-img/p2_phy_s2_q03_1.png}{140.5}{0}
\qfigdef{p2_phy_s2_q04_1}{JEE Advanced (2021 Series) - Test 05 (11 Apr 2021)-img/p2_phy_s2_q04_1.png}{144.9}{0}
\qfigdef{p2_phy_s3_q01_2}{JEE Advanced (2021 Series) - Test 05 (11 Apr 2021)-img/p2_phy_s3_q01_2.png}{108.6}{0}
\qfigdef{p2_phy_s3_q02_1}{JEE Advanced (2021 Series) - Test 05 (11 Apr 2021)-img/p2_phy_s3_q02_1.png}{96.0}{0}
\qfigdef{p2_phy_s3_q06_1}{JEE Advanced (2021 Series) - Test 05 (11 Apr 2021)-img/p2_phy_s3_q06_1.png}{92.7}{0}
\qfigdef{p2_chem_s1_q02_D}{JEE Advanced (2021 Series) - Test 05 (11 Apr 2021)-img/p2_chem_s1_q02_D.png}{56.9}{0}
\qfigdef{p2_chem_s1_q05_1}{JEE Advanced (2021 Series) - Test 05 (11 Apr 2021)-img/p2_chem_s1_q05_1.png}{152.3}{0}
\qfigdef{p2_chem_s1_q05_2}{JEE Advanced (2021 Series) - Test 05 (11 Apr 2021)-img/p2_chem_s1_q05_2.png}{215.1}{0}
\qfigdef{p2_chem_s1_q05_3}{JEE Advanced (2021 Series) - Test 05 (11 Apr 2021)-img/p2_chem_s1_q05_3.png}{214.9}{0}
\qfigdef{p2_chem_s2_q06_1}{JEE Advanced (2021 Series) - Test 05 (11 Apr 2021)-img/p2_chem_s2_q06_1.png}{175.2}{0}
\qfigdef{p2_chem_s3_q02_1}{JEE Advanced (2021 Series) - Test 05 (11 Apr 2021)-img/p2_chem_s3_q02_1.png}{407.4}{1}
\qfigdef{p2_chem_s3_q02_3}{JEE Advanced (2021 Series) - Test 05 (11 Apr 2021)-img/p2_chem_s3_q02_3.png}{426.4}{1}
\qfigdef{p2_chem_s3_q02_6}{JEE Advanced (2021 Series) - Test 05 (11 Apr 2021)-img/p2_chem_s3_q02_6.png}{125.5}{0}
\qfigdef{p2_chem_s3_q05_1}{JEE Advanced (2021 Series) - Test 05 (11 Apr 2021)-img/p2_chem_s3_q05_1.png}{52.2}{0}
\qfigdef{p2_chem_s3_q05_4}{JEE Advanced (2021 Series) - Test 05 (11 Apr 2021)-img/p2_chem_s3_q05_4.png}{75.9}{0}
\qfigdef{p2_chem_s3_q05_3}{JEE Advanced (2021 Series) - Test 05 (11 Apr 2021)-img/p2_chem_s3_q05_3.png}{28.4}{0}
\qfigdef{p2_chem_s3_q05_2}{JEE Advanced (2021 Series) - Test 05 (11 Apr 2021)-img/p2_chem_s3_q05_2.png}{82.9}{0}
\qfigdef{p2_chem_s3_q05_5}{JEE Advanced (2021 Series) - Test 05 (11 Apr 2021)-img/p2_chem_s3_q05_5.png}{86.2}{0}
\qfigdef{p2_chem_s3_q06_1}{JEE Advanced (2021 Series) - Test 05 (11 Apr 2021)-img/p2_chem_s3_q06_1.png}{116.3}{0}

\begin{document}
\begin{center}{\Large\bfseries JEE Advanced (2021 Series) --- Test 05 (11 Apr 2021)}\end{center}
\begin{multicols}{2}
\papersec{1}
\partsec{1}{PHYSICS}
\sect{Section-I(i)}
\qn{1} Four charges of $6\mu$C, $2\mu$C, $-12\mu$C and $4\mu$C are placed at the circumference of a circle. The circle is in $x$-$y$ plane and its centre at its origin. Locus of the points where electric potential is zero, is
\begin{qoptions}
\opt{A}{$x=y,\ z=0$}
\opt{B}{$x=0=z$}
\opt{C}{$x=0=y$}
\opt{D}{$x=z,\ y=z$}
\end{qoptions}
\qfullbegin
\qn{2} In the figure shown, the numbers indicate the values of resistance in ohm. The equivalent resistance of the combination across AB is :
\qfig{p1_phy_s1_q02_1}
\qfullsplit
\begin{qoptions}
\opt{A}{$\left(\dfrac{3+\sqrt{17}}{4}\right)$}
\opt{B}{$3+\sqrt{17}$}
\opt{C}{$\dfrac{3+\sqrt{17}}{2}$}
\opt{D}{$2(3+\sqrt{17})$}
\end{qoptions}
\qfullend
\qn{3} A sphere of radius $r$ is undergoing pure rolling with constant angular velocity on a fixed curved surface of radius $R$ with O as centre as shown. The ratio of the angle through which line PC rotates to the angle through which line OC rotates is :
\qfig{p1_phy_s1_q03_1}
\begin{qoptions}
\opt{A}{$\dfrac{R+r}{r}$}
\opt{B}{$\dfrac{R-r}{r}$}
\opt{C}{$\dfrac{R+r}{R}$}
\opt{D}{$\dfrac{R-r}{R}$}
\end{qoptions}
\qn{4} Figure shows combination of two semi infinite long wires kept on Z and Y-axis. It carries constant current $I$ as shown in figure. Calculate $\displaystyle\int_Q^P\vec B\cdot\vec{d\ell}$ on line QP. Coordinate of point P and Q are $(a,0,0)$ and $(0,0,a)$ respectively.
\qfig{p1_phy_s1_q04_1}
\begin{qoptions}
\opt{A}{$\mu_0I$}
\opt{B}{$\dfrac{\mu_0I}{2}$}
\opt{C}{$\dfrac{\mu_0I}{4}$}
\opt{D}{$\dfrac{\mu_0I}{8}$}
\end{qoptions}
\qn{5} A standing wave is produced in a vapour of an element of atomic mass $254$ (g mol$^{-1}$) at $400$ K. Nodes are found to be $6.77$ cm apart when frequency of source was $1000$ Hz. The vapour used for the experiment is ($R=8.31$ JK$^{-1}$ mol$^{-1}$)
\begin{qoptions}
\opt{A}{monatomic}
\opt{B}{diatomic}
\opt{C}{polyatomic}
\opt{D}{cannot be determined}
\end{qoptions}
\qn{6} The energy that should be added to an electron, to reduce its de-Broglie wavelength from $2\times10^{-9}$ m to $0.5\times10^{-9}$ m will be closest to :
\begin{qoptions}
\opt{A}{$1.1$ MeV}
\opt{B}{$40.56$ MeV}
\opt{C}{$0.56$ KeV}
\opt{D}{$5.6$ eV}
\end{qoptions}
\sect{Section-I(ii)}
\qn{7} In a series $L-C-R$ circuit, the voltage across resistance, capacitance and inductance are $30$V each at frequency $f=f_0$.
\begin{qoptions}
\opt{A}{If the inductor is short-circuited, the voltage across the capacitor will be $30\sqrt2$ V.}
\opt{B}{If the capacitor is short-circuited, the voltage drop across the inductor will be $\dfrac{30}{\sqrt2}$ V.}
\opt{C}{If the frequency is changed to $2f_0$, the ratio of reactance of the inductor to that of the capacitor is $4:1$.}
\opt{D}{If the frequency is changed to $2f_0$, the ratio of the reactance of the inductor to that of the capacitor is $1:4$.}
\end{qoptions}
\qn{8} A charged particle of charge $+q$ and mass $m$ is projected along $x$-axis with a speed $u$ and enters a magnetic field at point O.\\
The magnetic field is confined to a region\\
\hspace*{1em}$0\le x<0.8r,\ 0\le y\le R$\\
The gravitational field is confined to a region\\
\hspace*{1em}$0\le x\le(2r),\ R\le y\le\infty$\\
The electric field is confined to a region\\
\hspace*{1em}$(2r)\le x\le\infty,\ 0\le y\le\infty,\ R=\dfrac{2mu}{Bq}=2r$
\qfig{p1_phy_s2_q08_1}
\begin{qoptions}
\opt{A}{The charged particle will not enter into gravitational field}
\opt{B}{The charged particle will not enter into electric field}
\opt{C}{The charged particle will enter into both gravitational field and electric field one after the other}
\opt{D}{The charged particle will pass through the point C with same speed $u$}
\end{qoptions}
\qn{9} A part of a large circuit is shown. The values of resistance, inductance and capacitance are indicated here. Which of the statement(s) are correct?
\qfig{p1_phy_s2_q09_1}
\begin{qoptions}
\opt{A}{The charge on the capacitor is increasing with time}
\opt{B}{The energy of the inductor is increasing with time}
\opt{C}{The energy of inductor is decreasing with time}
\opt{D}{The energy of capacitor is decreasing with time}
\end{qoptions}
\qn{10} A potentiometer has a cell of emf $10$ V and internal resistance $2\Omega$. It is connected to a resistance box in which we can have any integral resistance from $1\Omega$ to $20\Omega$. The resistance of the potentiometer wire is $10\Omega$ and its length is $4.000$ m. The least count of the scale is $1$ mm. We wish to measure the emf of a secondary cell whose emf is $7$ V.
\begin{qoptions}
\opt{A}{The reading will be most accurate when we set the resistance box at $2\Omega$.}
\opt{B}{The reading will be most accurate when we set the resistance box at $3\Omega$.}
\opt{C}{If the resistance box is set at $1\Omega$ and we balance the secondary cell on the potentiometer wire, the null point will be at $3.640$ m.}
\opt{D}{If the resistance box is set at $8\Omega$ and we balance the secondary cell on the potentiometer wire, the null point will be at $2.000$ m.}
\end{qoptions}
\qn{11} A bowl of water is placed on one of the plates of a beam scale, which is balanced. One of our fingers is immersed into the water, such that the bowl is not touched. Will the scale remain balanced or not, if (i) not a drop of water flows out of the bowl; (ii) the bowl was full to the brim, and the water flows off the plate of the scale?
\begin{qoptions}
\opt{A}{In case (i) the scale remains balanced.}
\opt{B}{In case (i) the bowl of water weighs more.}
\opt{C}{In case (ii) the bowl of water weighs less.}
\opt{D}{In case (ii) the bowl of water weighs the same.}
\end{qoptions}
\qn{12} A small ball of radius $1$ cm is suspended in vacuum. It has a heater inside it which heats it at a rate of $72$ W. Take Stefan's constant to be $18/\pi\times10^{-8}$ W/m$^2$K$^4$. We have two possible paints with which we can paint the ball. One of them has an emissivity of $0.64$ and another paint is black.
\begin{qoptions}
\opt{A}{If the ball is painted with black paint, its surface temperature would be $1000$ K.}
\opt{B}{If the ball is painted with paint of emissivity $0.64$, its surface temperature would be $800$ K approx.}
\opt{C}{If the ball is half painted with one and other half is painted with other paint, the black part will be hotter.}
\opt{D}{If the ball is painted with black paint, the peak emissive power is corresponding to wavelength of $2880$ nm.}
\end{qoptions}
\sect{Section-II}
\qn{1} A point source of light lies along axis, outside a long glass cylinder as shown. The refractive index of glass is $\mu=1.5$. The distance of the source from the face A is $R$, where $R$ is the radius of the cylinder. If the power of the source is $50\sqrt2$ W, what is the power (in W) of the beam coming from the other flat face? Neglect all partial reflections.
\qfig{p1_phy_s3_q01_1}
\qn{2} One mole of ideal monoatomic gas is taken along a cyclic process as shown in the figure. Process $1\to2$ shown is $1/4^{\text{th}}$ part of a circle as shown by dotted lines, process $2\to3$ is isochoric while $3\to1$ is isobaric. Calculate efficiency in percentage of the above cycle.
\qfig{p1_phy_s3_q02_1}
\qn{3} A bumper is mounted to a vertical rod, hinged at its bottom end, at a distance of $d=10$ cm from the hinge. There is a small bead at rest on the bumper. The rod is started to vibrate and executes SHM about its original position with a small angular amplitude of $\theta_0=1^\circ$, as shown in the figure. What should the minimum frequency (in Hz) of the vibration be in order that the bead flies off? (Friction is negligible.) Take $g=10$ m/s$^2$ and $\pi^2=10$
\qfig{p1_phy_s3_q03_1}
\qn{4} The radioactive carbon-14 isotope is constantly produced in the atmosphere by the cosmic rays. Nonetheless, the amount of 14C on the Earth is constant, because at a height of 5-9 km, as a result of the reaction $n+{}^{14}N\to{}^{14}C+p$, an average of $18000$; 14C isotopes are produced in each second on each square metre -- projected to the surface of the Earth. The half-life of the isotope is $6930$ years. Estimate how many tons of 14C are present on the Earth. ($R_{\text{earth}}=6400$ km)
\qn{5} Three weightless pulleys with radii $R$, $2R$ and $3R$ are concentrically fastened together to make a single triple pulley system and mounted on an axis on which the triple pulley can rotate without friction (figure). The threads are light and they do not slip on the pulley. We suspended masses $10m$, $m$ and $2m$ as shown. What is the acceleration (in m/s$^2$) of mass $2m$.
\qfig{p1_phy_s3_q05_1}
\qn{6} A $1$ m long tube is open at one end and closed at the other end. Initially it contains a liquid column of $55$ cm which traps $25$ cm column of air when held horizontal. Now the tube is held vertical with open end at bottom. Now what will be the equilibrium height of liquid column (in cm), after fall of some liquid, if some liquid falls? (Neglect surface tension, assume temperature of the gas remains constant). $P_{\text{atm}}=10^5$ Pa, $\rho_{\text{liquid}}=10^2$ kg/m$^3$, $g=10$ m/s$^2$.
\qfig{p1_phy_s3_q06_1}
\partsec{2}{CHEMISTRY}
\sect{Section-I(i)}
\qn{1} $\mathrm{Cu^{2+}(aq.)+Zn(s)\longrightarrow Cu(s)+Zn^{2+}(aq.)}$ ; $E^0=1.1$ V
\[\mathrm{Zn^{2+}(aq.)+2e^-\longrightarrow Zn(s)}\ ;\ E^0=-0.76\text{ V}\]
\[\begin{gathered}\mathrm{Cu^{2+}(aq.)+4NH_3(aq.)\longrightarrow[Cu(NH_3)_4]^{2+}(aq.)}\\ K_{\mathrm{eq}}=10^{20}\end{gathered}\]
If $\dfrac{2.303RT}{F}=0.06$ then the $E^0$ of following reaction is -
\[\mathrm{Cu(s)+4NH_3(aq.)\longrightarrow[Cu(NH_3)_4]^{2+}(aq.)+2e^-}\]
\begin{qoptions}
\opt{A}{$0.34$ V}
\opt{B}{$0.76$ V}
\opt{C}{$0.26$ V}
\opt{D}{$0.14$ V}
\end{qoptions}
\qfullbegin
\qn{2} \qfigi{p1_chem_s1_q02_1}
\begin{qoptions}
\opt{A}{\qfigi{p1_chem_s1_q02_2}}
\opt{B}{\qfigi{p1_chem_s1_q02_3}}
\opt{C}{\qfigi{p1_chem_s1_q02_C}}
\opt{D}{\qfigi{p1_chem_s1_q02_D}}
\end{qoptions}
\qfullend
\qn{3} Which of the following having highest density among the following?
\begin{qoptions}
\opt{A}{Zn}
\opt{B}{Co}
\opt{C}{V}
\opt{D}{Cr}
\end{qoptions}
\qn{4} With the help of following phase diagram, select the only \textbf{INCORRECT} statement -
\qfig{p1_chem_s1_q04_1}
\begin{qoptions}
\opt{A}{Mole fraction of B in vapour at pressure $P_1$ for solution $Z_2$ is $0.5$}
\opt{B}{Mole fraction of A in liquid at pressure $P_3$ for solution $Z_2$ is $0.4$}
\opt{C}{Mole fraction of B in vapour at pressure $P_1$ for solution $Z_2$ is $0.75$}
\opt{D}{Mole fraction of A in liquid at pressure $P_2$ for solution $Z_1$ is $0.4$}
\end{qoptions}
\qn{5} \[\begin{gathered}\mathrm{CH_3{-}CH_2{-}C{\equiv}CH{-}CH_2{-}\overset{\overset{O}{\|}}{C}{-}Cl}\ \xrightarrow[\mathrm{BaSO_4}]{\mathrm{H_2,Pd,(excess)}}\ \mathrm{(X)}\\ \mathrm{(X)}\ \xrightarrow[\mathrm{HOH}]{\mathrm{CH_3MgBr}}\ \mathrm{Y+Z}\end{gathered}\]
Products Y and Z are
\begin{qoptions}
\opt{A}{Enantiomeric pair of \qfigi{p1_chem_s1_q05_1}}
\opt{B}{Diastereomeric pair of \qfigi{p1_chem_s1_q05_2}}
\opt{C}{Enantiomeric pair of \qfigi{p1_chem_s1_q05_3}}
\opt{D}{Product (Y) and (Z) are not stereoisomers}
\end{qoptions}
\qn{6} What is the overall complex dissociation equilibrium constant for the $\mathrm{[Cu(NH_3)_4]^{2+}}$ ion given that $\beta_4$ for the complex is $2.0\times10^{13}$.
\begin{qoptions}
\opt{A}{$2.0\times10^{13}$}
\opt{B}{$5\times10^{-14}$}
\opt{C}{$2\times10^{-13}$}
\opt{D}{data not sufficient}
\end{qoptions}
\sect{Section-I(ii)}
\qn{7} Radionuclide `P' decays into `Q' and `Q' decays into `R' as :
\[\mathrm{P}\xrightarrow{t_{1/2}=1000\text{ hrs}}\mathrm{Q}\xrightarrow{t_{1/2}=10\text{ min}}\mathrm{R}\]
Number of nuclei of Q and R after 1000 hr, if we start with $2$ mol P, are. $[N_A=6\times10^{23}]$
\begin{qoptions}
\opt{A}{$N_Q=6\times10^{23}$}
\opt{B}{$N_Q=10^{20}$}
\opt{C}{$N_R=6\times10^{23}$}
\opt{D}{$N_R=10^{23}$}
\end{qoptions}
\qn{8} Which of the following compounds can show positive test with tollen's reagent?
\begin{qoptions}
\opt{A}{$\mathrm{H{-}\overset{\overset{O}{\|}}{C}{-}OH}$}
\opt{B}{\qfigi{p1_chem_s2_q08_B}}
\opt{C}{\qfigi{p1_chem_s2_q08_C}}
\opt{D}{\qfigi{p1_chem_s2_q08_D}}
\end{qoptions}
\qn{9} Identify disproportionation reaction(s).
\begin{qoptions}
\opt{A}{$\mathrm{XeF_4\xrightarrow{H_2O}}$}
\opt{B}{$\mathrm{XeF_6\xrightarrow[\text{excess of }OH^-]{H_2O}}$}
\opt{C}{$\mathrm{Cl_2\xrightarrow{NaOH}}$}
\opt{D}{$\mathrm{NO_2\xrightarrow{H_2O}}$}
\end{qoptions}
\qn{10} A $1$ kg real gas is liquified at temperature $300$ K as shown in diagram. Select the correct statement-
\qfig{p1_chem_s2_q10_1}
\begin{qoptions}
\opt{A}{$2$ atm is vapour pressure of liquid at $300$ K}
\opt{B}{Density of gas when $V$ equals to $10$ litre, is $0.1$ gm/ml}
\opt{C}{Density of gas when $V$ equals to $10$ litre, is $0.05$ gm/ml}
\opt{D}{Density of liquid when $V$ equals to $10$ litre, is $1$ gm/ml}
\end{qoptions}
\qn{11} Which of the following alchols can change the coloure of acidic dicromate solution :
\begin{qoptions}
\opt{A}{\qfigi{p1_chem_s2_q11_A}}
\opt{B}{\qfigi{p1_chem_s2_q11_B}}
\opt{C}{\qfigi{p1_chem_s2_q11_C}}
\opt{D}{\qfigi{p1_chem_s2_q11_D}}
\end{qoptions}
\qn{12} Which of the following statement(s) is/are incorrect?
\begin{qoptions}
\opt{A}{Hydrometallurgy is used to recover Cu from Cu scraps}
\opt{B}{Hydrometallurgy is used for extraction of Ag from their native ore}
\opt{C}{$\mathrm{Na_2CO_3}$ is used as leaching agent in Hall's Heroult method}
\opt{D}{Cupellation is used to remove Pb from impure Ag/Au}
\end{qoptions}
\sect{Section-II}
\qn{1} A perfect diatomic gas at $500$ K undergoes following steps of process-\\
(a) An isochoric change of state from ($0.5$ atm, $80$ litre) to ($1$ atm, $80$ litre)\\
(b) An isobaric change from ($1$ atm, $80$ litre) to ($1$ atm, $20$ litre).\\
If all processes are irreversible then calculate the magnitude of entropy change (in J/K) of surrounding for entire process.\\
{[}Use: $\ln2=0.7$, $R=0.08$ atm-litre/K-mole, $1$ lit-atm $=100$ J.]
\qn{2} Formation of (P) is known as Tollen's reaction. Identify total number of stereoisomers for `Q'
\[\mathrm{CH_3CH_2CHO+3HCHO}\xrightarrow[\Delta]{\mathrm{NaOH}}\mathrm{P}\xrightarrow[\mathrm{H^+}]{\mathrm{CH_3CHO\ (1\ eq.)}}\mathrm{Q}\]
\qn{3} Find the number of chemical species having maximum number of atoms may lye in one plane is $5$.\\
$\mathrm{XeF_4}$, $\mathrm{PCl_5}$, $\mathrm{XeF_5^+}$, $\mathrm{B(OH)_3}$, $\mathrm{[Ni(CN)_4]^{2-}}$, $\mathrm{SF_6}$, $\mathrm{[Zn(CN)_4]^{2-}}$
\qn{4} Packing fraction of two dimensional unit cell shown in figure is
\qfig{p1_chem_s3_q04_1}
[Take $\pi=3.2$, $\sqrt2=1.4$]
\qn{5} How many reactions are written with their correct product?\par\medskip
\noindent(a)~\qfigi[0.8]{p1_chem_s3_q05_1}\par\medskip
\noindent(b)~\qfigi[0.8]{p1_chem_s3_q05_3}\par\medskip
\noindent(c)~\qfigi[0.8]{p1_chem_s3_q05_5}\par\medskip
\noindent(d)~\qfigi[0.8]{p1_chem_s3_q05_7}\par\medskip
\noindent(e)~\qfigi[0.8]{p1_chem_s3_q05_9}
\qn{6} Find the number of elements from the given set have greater tendency to convert in $\mathrm{M^{+2}}$ to $\mathrm{M^{+3}}$ than Cr.\\
Ti, Fe, V, Mn, Co
\columnbreak\partsec{3}{MATHEMATICS}
\sect{Section-I(i)}
\qn{1} If ABCD is regular tetrahedron with length of any edge `$\ell$', then the minimum distance of any vertex from the opposite face is -
\begin{qoptions}
\opt{A}{$\sqrt{\dfrac23}\,\ell$}
\opt{B}{$\sqrt{\dfrac32}\,\ell$}
\opt{C}{$\dfrac23\ell^2$}
\opt{D}{$\dfrac{1}{\sqrt3}\,\ell$}
\end{qoptions}
\qn{2} If $f(x)$ is continuous and differentiable over $[-2,5]$ and $-4\le f'(x)\le3$ for all $x$ in $(-2,5)$, then the greatest possible value of $f(5)-f(-2)$ is -
\begin{qoptions}
\opt{A}{$7$}
\opt{B}{$9$}
\opt{C}{$15$}
\opt{D}{$21$}
\end{qoptions}
\qn{3} Area bounded by $f(x)=\min(\{x-2\},\{3-x\})$, $x=4$, $x=-4$ and $x$-axis is equal to (where $\{.\}$ denotes fractional part)
\begin{qoptions}
\opt{A}{$1$}
\opt{B}{$3/2$}
\opt{C}{$2$}
\opt{D}{$4$}
\end{qoptions}
\qn{4} Circles $C_1: x^2+y^2=9$ and $C_2: x^2+y^2-4x=0$. Common chord $PQ$ of $C_1$ \& $C_2$ intersects the direct common tangent at $R$. If $RP+RQ=\sqrt{\dfrac ab}$, where $a$ \& $b$ are relatively prime. Then $(a+b)$ equals-
\begin{qoptions}
\opt{A}{$65$}
\opt{B}{$63$}
\opt{C}{$71$}
\opt{D}{$79$}
\end{qoptions}
\qn{5} Let $\vec a=\hat i+3\hat j+5\hat k$, $\vec b=3\hat i+5\hat j+7\hat k$, $\vec c=\hat i+\hat j-\hat k$. If the volume of parallelopiped formed by the vectors $2\vec a+\vec b+\vec c$, $\vec a-\vec b+2\vec c$ and $\vec a+3\vec b+\vec c$ is $v$, then
\begin{qoptions}
\opt{A}{$v\in[8,25]$}
\opt{B}{$v\in(25,50)$}
\opt{C}{$v\in[50,75]$}
\opt{D}{$v>75$}
\end{qoptions}
\qn{6} The equation of the parabola whose focus is the point $(0,0)$ and the tangent at the vertex is $x-y+1=0$, is
\begin{qoptions}
\opt{A}{$x^2+y^2-2xy-4x+4y-4=0$}
\opt{B}{$x^2+y^2-2xy+4x-4y-4=0$}
\opt{C}{$x^2+y^2+2xy-4x+4y-4=0$}
\opt{D}{$x^2+y^2+2xy-4x-4y+4=0$}
\end{qoptions}
\sect{Section-I(ii)}
\qn{7} If $\phi(x)=3f\!\left(\dfrac{x^2}{3}\right)+f(3-x^2)\ \forall\ x\in(-3,4)$ where $f''(x)>0\ \forall\ x\in(-3,4)$, then $\phi(x)$ is -
\begin{qoptions}
\opt{A}{increasing in $\left(\dfrac32,4\right)$}
\opt{B}{decreasing in $\left(-3,-\dfrac32\right)$}
\opt{C}{increasing in $\left(-\dfrac32,0\right)$}
\opt{D}{decreasing in $\left(0,\dfrac32\right)$}
\end{qoptions}
\qn{8} Consider two sets $A=\{1,2,3\}$ \& $B=\{1,2,3,4,5\}$. A number is selected each from set $A$ \& set $B$. Let $p$ denotes the probability that the selected number from set $A$ is not more than the selected number from set-$B$ \& $q$ denotes the probability that the number from set-$B$ is not more than the number from set $A$, then -
\begin{qoptions}
\opt{A}{$p=2q$}
\opt{B}{$2p=3q$}
\opt{C}{$p+q=\dfrac45$}
\opt{D}{$p+q=\dfrac65$}
\end{qoptions}
\qn{9} $P(x)$ is a least degree polynomial such that $(P(x)-1)$ is divisible by $(x-1)^2$ \& $(P(x)-3)$ is divisible by $(x+1)^2$, then
\begin{qoptions}
\opt{A}{graph of $y=P(x)$ is symmetric about origin}
\opt{B}{$y=P(x)$ has two points of extrema}
\opt{C}{$\displaystyle\int_{-4}^{\lambda}P(x)\,dx=0$ for exactly one value of $\lambda$.}
\opt{D}{$\displaystyle\int_{-4}^{\lambda}P(x)\,dx=0$ for exactly two values of $\lambda$.}
\end{qoptions}
\qn{10} If $z_1,z_2,z_3\in\mathbb C$ and satisfying $|2z_1-1|=4$, $|2z_2-5|=1$ and $|z_3-3|=1$ respectively, then
\begin{qoptions}
\opt{A}{least value of $|z_1+z_2+z_3|$ is $\dfrac32$}
\opt{B}{least value of $|z_1+z_2+z_3|$ is $\dfrac52$}
\opt{C}{greatest value of $|z_1+z_2+z_3|$ is $\dfrac{19}{2}$}
\opt{D}{greatest value of $|z_1+z_2+z_3|$ is $\dfrac{11}{2}$}
\end{qoptions}
\qn{11} Equation of curve passing through $(3,4)$ and satisfying the differential equation, $y\left(\dfrac{dy}{dx}\right)^2+(x-y)\dfrac{dy}{dx}-x=0$ can be
\begin{qoptions}
\opt{A}{$x^2+y^2=25$}
\opt{B}{$x^2-y^2+7=0$}
\opt{C}{$xy=12$}
\opt{D}{$x-y+1=0$}
\end{qoptions}
\qn{12} Let
\[f(x)=\begin{cases}e^{\tan^{-1}(x-1)}+(x-1)^{2017}\sin^2\dfrac{1}{(x-1)} & \text{if } x\neq1\\[4pt] 1 & \text{if } x=1\end{cases}\]
and
\[\ell=\lim_{n\to\infty}\left(n\sum_{k=1}^{10}f\!\left(1+\dfrac kn\right)-10n\right)\]
\begin{qoptions}
\opt{A}{$f(x)$ is continuous but non-derivable at $x=1$}
\opt{B}{$f(x)$ is differentiable at $x=1$}
\opt{C}{$55f'(1)-\ell=0$}
\opt{D}{$\ell$ is divisible by $11$.}
\end{qoptions}
\sect{Section-II}
\qn{1} Let $\vec x,\vec y$ are non zero and non collinear vectors such that
\[(a\alpha^2+b\alpha+c)\vec x+(a\beta^2+b\beta+c)\vec y+(a\gamma^2+b\gamma+c)(\vec x\times\vec y)=0,\]
then $(a+b+c)$ is (where $a,b,c\in\mathbb R$)
\qn{2} Let $P(x)$ be a polynomial with degree $4$ such that $\displaystyle\lim_{x\to0}\left(\dfrac{P(x)}{x}-12\right)=0$. The graph of $P(x)$ passes through $A(1,25)$ \& $B(-1,-7)$ \& has extremum at point $B$. If the equation $|P(x)|=k$ has exactly $3$ roots, then $k$ is
\qn{3} Consider $I_1=\displaystyle\int_0^{\pi/2}(\cos x)^{\sqrt{17}+3}\,dx$ and $I_2=\displaystyle\int_0^{\pi/2}(\cos x)^{\sqrt{17}+5}\,dx$, then $\dfrac{I_1}{I_2}$ is equal to (Use $\sqrt{17}=4.12$)-
\qn{4} Slope of normal to the ellipse at a point $P$ is $\dfrac34$ and eccentricity of ellipse is $\dfrac13$. If this normal makes acute angle `$\beta$' with its focal chord through $P$, then $\sin\beta$ is
\qn{5} Let $A$ is a four order diagonal matrix whose entries are complex numbers such that $A^4=I_4$. If trace of $A$ is zero, then number of such matrices is\\
(Here $I_4$ is four order unit matrix)
\qn{6} Consider a sequence of distinct terms $\{a_n\}$ with $a_1=2$ \& $a_n=\dfrac{a_{n-1}^2}{a_{n-2}}$ for $\forall\ n\ge3$. If $a_5\le162$ and $a_2$, $a_5$ are positive integers than sum of possible value(s) of $a_5$ is
\papersec{2}
\partsec{1}{PHYSICS}
\sect{Section-I}
\qn{1} In the given circuit switch is closed at $t=0$,
\qfig{p2_phy_s1_q01_1}
\begin{qoptions}
\opt{A}{Current in the inductor when the circuit reaches the steady state is $4$ A.}
\opt{B}{The net change in flux in the inductor is $1.5$ Wb.}
\opt{C}{The time constant of the circuit after closing S is $555.55$ s.}
\opt{D}{The charge stored in the capacitor in the steady state is $1.2$ mC.}
\end{qoptions}
\qn{2} In the Bohr model of the atom (Here $n$ is the principal quantum number)
\begin{qoptions}
\opt{A}{Frequency of revolution of the electron is inversely proportional to $n^3$}
\opt{B}{The modulus of potential energy of the electron in an orbit is greater than its kinetic energy.}
\opt{C}{Magnetic induction at the nucleus produced due to the motion of electron motion in the $n^{\text{th}}$ orbit is proportional to $n^5$.}
\opt{D}{Magnetic moment produced due to the motion of electron motion in the $n^{\text{th}}$ orbit is proportional to $n$.}
\end{qoptions}
\qn{3} Two light bodies are kept floating on the surface of water very close to each other as shown below. Water surface takes different shapes depending on the surface nature of the bodies. Identify true situations from the following arrangement in the view of surface tension : Assume complete wetting/$\theta=90^\circ$/ no wetting as the case may be.
\begin{qoptions}
\opt{A}{\qfigi{p2_phy_s1_q03_A} both the bodies get pulled towards.}
\opt{B}{\qfigi{p2_phy_s1_q03_B} both the bodies get pulled towards.}
\opt{C}{\qfigi{p2_phy_s1_q03_C} both the bodies get repelled}
\opt{D}{\qfigi{p2_phy_s1_q03_D} neither repulsion nor attraction}
\end{qoptions}
\qn{4} $n$ moles of an ideal gas undergo an isothermal process at temperature $T$. $P$-$V$ graph of the process is as shown in the figure. A point $A(V_1,P_1)$ is located on the $P$-$V$ curve. Tangent at point $A$, cuts the $V$-axis at point $B$. $AO$ is the line joining the point $A$ to the origin $O$ of $PV$ diagram. $BC$ is perpendicular to $OB$. Then
\qfig{p2_phy_s1_q04_1}
\begin{qoptions}
\opt{A}{The coordinates of point $C$ are $\left(2V_1,\dfrac{P_1}{2}\right)$}
\opt{B}{The coordinates of point $C$ are $\left(\dfrac32V_1,\dfrac23P_1\right)$}
\opt{C}{Area of the triangle $AOB$ is equal to $nRT$}
\opt{D}{Area of the triangle $AOB$ is equal to $\dfrac{nRT}{2}$}
\end{qoptions}
\qn{5} Given figure shows how the temperature of water and ice mixed with each other at different temperatures, changes with the amount of heat ($Q$) evolved by the water or absorbed by the ice. The scale corresponding to heat ($Q$) is arbitrary. Which of the following is/are conclusions drawn from the graph ?
\qfig{p2_phy_s1_q05_1}
\begin{qoptions}
\opt{A}{The mass of water is double that of ice.}
\opt{B}{The equilibrium temperature will be $0^\circ$C with three-fourth of the ice melted.}
\opt{C}{The equilibrium temperature will be $0^\circ$C with one fourth of the water frozen.}
\opt{D}{The equilibrium temperature will be $10^\circ$C.}
\end{qoptions}
\qn{6} In the circuit in Figure, the ``horizontal'' battery has a voltage of $1$ V, three of the four capacitors have the same capacity and the last twice as big. What can be the voltage of a second, ``vertical'' battery of voltage $U_x$, so that at least one capacitor in this circuit remained uncharged? Before battery connections were made capacitors are not charged.
\qfig{p2_phy_s1_q06_1}
\begin{qoptions}
\opt{A}{$1$ V}
\opt{B}{$1.5$ V}
\opt{C}{$2/3$ V}
\opt{D}{$2$ V}
\end{qoptions}
\sect{Section-II}
\qn{1} There is an isolated planet having mass $2M$ and radius $2R$, where $M$ and $R$ are the mass and radius of the earth. A simple pendulum having mass $m$ and length $2R$ is made to small oscillations on the planet. Find the time period of SHM of pendulum (in second).\\
(take $\pi=3.00$, $g=10$ m/s$^2$, $R=6400$ km, $\sqrt2=1.41$)
\qn{2} A uniform rod of mass $m$ and length $L=2$ m is clamped at mid point, vibrates in the first overtone, with the frequency of $5000$ Hz. Find the density of the rod (in kg/m$^3$). ($Y$ is the young modulus of the material of the rod $=2\times10^{11}$ N/m$^2$).
\qn{3} A slender rod of mass $M$ and length $L=1$ m is lying on the frictionless surface. A particle of mass $\dfrac M8$ moving perpendicular to the length of the rod, with speed $v$ and strikes the rod at a distance $\dfrac L4$ from the centre of mass of the rod. Find the $y$ coordinate (in m) of the end $B$ in time $t=\left(\dfrac{4\pi L}{v}\right)$ (Assume that after collision particle comes to rest.)
\qfig{p2_phy_s2_q03_1}
\qn{4} In free space on a circle with radius $R_0$, four point masses $m$ are located at the vertices of the inscribed square, two of them carry charge $+q$, and the other two charge $-q$. At the initial moment all these particles have same speed in clockwise direction as shown. It is known that during the motion the minimum distance from any of the point masses to the centre $O$ of the initial circle is $R_1$ ($R_1<R_0$). Consider that at any time, the charges are at the vertices of square centered at point $O$. The action of gravitational forces can be neglected. Determine the speed (in m/s) of any particle at a position having a distance $R_1$ from the centre of the circle. Take $m=2\sqrt2-1$ gm, $q=1\,\mu$C, $R_0=2R_1=1$ m.
\qfig{p2_phy_s2_q04_1}
\qn{5} A point source of light emits a power of $100$ W of monochromatic light with a wavelength of $248$ nm. This source is at the focus of a convex lens of focal length $10$ cm and aperture diameter $2$ cm. On the other side of the lens we have a small metallic plate of radius $1$ mm kept perpendicular to the optic axis such that its centre lies on the optic axis. The metallic plate has work function of $1.24$ eV and a photoelectric efficiency of $0.1\%$. What is the initial photocurrent (in $\mu$A) from the plate after the source is switched on ?
\qn{6} If the beat frequency can be measured up to an accuracy of $0.1\%$, find the possible error (in cm/s) in measurement of speed of a vehicle moving with a approximate speed of $30$ m/s towards a stationary source. The detector is located behind the source and is stationary. The source, detector and the moving vehicle are always on the same straight line. The speed of sound in air is known to be $330$ m/s. The frequency is accurately known.
\sect{Section-III}
\qn{1} A solid block of mass $2$ kg is resting inside a cube as shown in figure. The cube is moving with a velocity $\vec v=5\hat i+2t\hat j$ m/s. If the coefficient of friction between the surface of cube and block is $0.2$. Then what is the force of friction between the block and cube (in N) ?
\qfig{p2_phy_s3_q01_2}
\qn{2} A uniformly charged rod of charge per unit length $\lambda$, is rotated about an axis as shown in the figure with a constant angular velocity $\omega=\dfrac{6}{(\lambda\ell^3)}$. Find the magnetic moment of the system.
\qfig{p2_phy_s3_q02_1}
\qn{3} A particle of mass `$m$' having collided with a stationary particle of mass `$M$' deviated by an angle $\dfrac\pi2$ rad where as the particle of mass `$M$' moving at an angle of $30^\circ$ to the direction of the initial motion of the particle `$m$'. If after the collision there is $40\%$ loss in kinetic energy of the two particle system, find $\dfrac Mm$.
\qn{4} Assume that the refractive index of atmosphere of the earth is uniform and equal to $1.003$. Assume that the atmosphere extends to height of $80$ km above the surface of the earth. The radius of the earth is $6400$ km. Due to the refraction, the sun sets after it has gone just below the horizon. Find the angular position of sun below the horizon (in degrees) when we see the sun setting. Round off to nearest integer.
\qn{5} A body is executing simple harmonic motion. At a displacement $x$ from mean position, its potential energy is $E_1=1$ J and at a displacement $y$ from mean position, its potential energy is $E_2=4$ J. The potential energy $E$ at a displacement $(x+y)$ from mean position is \underline{\hspace{2cm}} J
\qn{6} A transparent sphere of radius $R$ has a cavity of radius $\dfrac R2$ as shown in figure. Find the refractive index of the sphere if a narrow parallel beam of light falling on left surface focuses at point $P$. If this refractive index is $\mu=\dfrac{3+\sqrt X}{2}$, find $x$.
\qfig{p2_phy_s3_q06_1}
\partsec{2}{CHEMISTRY}
\sect{Section-I}
\qn{1} Which of the following order is/are correct for hydrides of 15th group ($\mathrm{EH_3}$) where $E=N/P/As/Sb$
\begin{qoptions}
\opt{A}{Highest boiling point $\mathrm{NH_3}$}
\opt{B}{Highest melting point $\mathrm{SbH_3}$}
\opt{C}{Highest E--H distance $\mathrm{SbH_3}$}
\opt{D}{Highest $\Delta H_{\mathrm{disso.}}$ E--H $\mathrm{NH_3}$}
\end{qoptions}
\qn{2} Which of the following reaction gives aldehyde or ketone as final product?
\begin{qoptions}
\opt{A}{$\mathrm{(CH_3CO_2)_2Ca\xrightarrow{\Delta}}$}
\opt{B}{$\mathrm{Ph{-}\overset{\overset{O}{\|}}{C}{-}Cl\xrightarrow[\text{(Quinoline)}]{H_2\,(Pd-BaSO_4)}}$}
\opt{C}{$\mathrm{H_3C{-}CH_2{-}C{\equiv}CH\xrightarrow[H_3O^+]{HgSO_4}}$}
\opt{D}{\qfigi{p2_chem_s1_q02_D}$\ \xrightarrow{\text{Hot }\mathrm{KMnO_4}}$}
\end{qoptions}
\qn{3} $1$ litre aq. solution containing sufficient sodium salt of 2,3 dimethyl butan-1,4dioicacid is electrolysed in an electrolytic cell using $9.65$ mA current for $2\times10^6$ sec using inert Pt electrode with $100\%$ current efficiency. Select the correct option(s) -
\begin{qoptions}
\opt{A}{Mass of hydrocarbon formed at anode is $5.6$ gm}
\opt{B}{Mass of hydrocarbon formed at anode is $11.2$ gm}
\opt{C}{Gas evolved at cathode occupies $2.27$ litre at STP}
\opt{D}{Gas evolved at cathode occupies $1.135$ litre at STP}
\end{qoptions}
\qn{4} Which of the following statement(s) is/are CORRECT?
\begin{qoptions}
\opt{A}{D-penicillamine is used to remove excess of Cu through chelate therapy}
\opt{B}{Desferrioxamine is used to remove excess of Fe through chelate therapy}
\opt{C}{EDTA is used in the treatment of lead poisoning}
\opt{D}{cis-platin and related compound effectively inhibit the growth of tumour}
\end{qoptions}
\columnbreak\qn{5} \qfig{p2_chem_s1_q05_1}
correct statement(s) about above reactions and products is/are
\begin{qoptions}
\opt{A}{Major product X is\\ \qfigi[0.8]{p2_chem_s1_q05_2}}
\opt{B}{Major product X is\\ \qfigi[0.8]{p2_chem_s1_q05_3}}
\opt{C}{In HOH, $\overline{\mathrm O}$ is highly solvated through hydrogen bond.}
\opt{D}{In DMSO, no hydrogen bond is formed}
\end{qoptions}
\qn{6} Select the incorrect statement(s) :
\begin{qoptions}
\opt{A}{Stoichiometric defects always lowers density of solid.}
\opt{B}{Defect only appears in ionic solids.}
\opt{C}{Alkaline earth flouride may show anion Frenkel defect.}
\opt{D}{When AgCl is dopped with $\mathrm{Ag_2S}$, then anion vacancy created.}
\end{qoptions}
\sect{Section-II}
\qn{1} What is the sum of the square of the magnetic moments (spin only) in the following complex.\\
$\mathrm{[Ni(NH_3)_2Cl_2]}$, $\mathrm{[RhF_6]^{3-}}$, $\mathrm{[Cu(NH_3)_4]^{2+}}$, $\mathrm{[PtCl_4]^{2-}}$, $\mathrm{[Ag(CN)_2]^-}$, $\mathrm{[MnCl_4]^{2-}}$, $\mathrm{Na_2[Fe(H_2O)_5NO]}$
\qn{2} Identify total number of reagent sequence are used for following conversion :
\[\mathrm{CH_3{-}Cl\longrightarrow CH_3{-}CH_3}\]
\begin{enumerate}[label=(\alph*),leftmargin=*,itemsep=0pt]
\item Na / Dry ether
\item NaCN / $\mathrm{H_3O^+}$ / NaOH / Electrolysis
\item $\mathrm{LiAlH_4}$
\item MeMgCl / Dry ether
\item Mg--Dry ether / Me -- I
\item Mg / Dry ether / $\mathrm{CO_2}$ / $\mathrm{LiAlH_4}$ / Red P + HI
\item Aq. NaOH
\end{enumerate}
\qn{3} When an electron jumps from $n^{\text{th}}$ orbit to the $1^{\text{st}}$ orbit in H-atoms in one or more than one transitions, the the longest wavelength of the emitted photons is $1872$ nm. The value of `$n$' is (Given : $\dfrac1R=91$ nm)
\qn{4} Find the number of reaction in which $\mathrm{CrO_4^{2-}}$ or $\mathrm{Cr_2O_7^{2-}}$ does not act as oxidising agent.
\[\mathrm{CrO_4^{2-}\xrightarrow{HNO_3}},\quad \mathrm{KI\xrightarrow{Cr_2O_7^{2-}/H^+}}\]
\[\mathrm{Fe^{+2}\xrightarrow{Cr_2O_7^{2-}/H^+}},\quad \mathrm{H_2O_2\xrightarrow{CrO_4^{2-}/OH^-}}\]
\[\mathrm{Sn^{+2}\xrightarrow{Cr_2O_7^{2-}/H^+}}\]
\qn{5} How many of the following can be used to convert ethanol into ethanal only\\
\mbox{(a) $\mathrm{H^{\oplus}KMnO_4/\Delta}$} \quad \mbox{(b) Cu/$\Delta$} \quad \mbox{(c) NBS}\\
\mbox{(d) DMSO + $\mathrm{NaHCO_3}$ /$\Delta$} \quad \mbox{(e) PCC}\\
\mbox{(f) $\mathrm{H^{\oplus}K_2Cr_2O_7}$ /$\Delta$} \quad \mbox{(g) $\mathrm{HIO_4}$ /$\Delta$} \quad \mbox{(h) Schiff reagent}
\qn{6} In $10$ ml, $0.2$M $\mathrm{CH_3COOH}$ ($K_a=10^{-5}$M) aq. solution, $0.1$M $\mathrm{NH_4OH}$ ($K_b=10^{-5}$M) is gradually added and following conductance graph is observed.
\qfig{p2_chem_s2_q06_1}
If pH of $\mathrm{CH_3COOH}$ solution initially is $\mathrm{pH_1}$ \& on adding $10$ ml base it is $\mathrm{pH_2}$, calculate ($\mathrm{pH_2-pH_1}$). [Given : $\log2=0.3$]
\sect{Section-III}
\qn{1} Which of the following interhalogen on hydrolysis does not form oxyacid of their underliend atom.\\
$\underline{\mathrm I}\mathrm F_7$, $\mathrm{Cl}\underline{\mathrm{Br}}$, $\underline{\mathrm{Cl}}\mathrm F_5$, $\underline{\mathrm{Br}}\mathrm F_3$
\qfullbegin
\qn{2} Identify total number of reactions which give Aromatic product ?
\qfig{p2_chem_s3_q02_1}
\qfullsplit
\qfig{p2_chem_s3_q02_3}
\qfig{p2_chem_s3_q02_6}
\qfullend
\qn{3} `A' decomposes to give `X' with a rate constant of $6.93\times10^{-2}$ min$^{-1}$. The reactant `A' in the presence of catalyst gives `Y' by a second order mechanism with a rate constant of $0.02$ M$^{-1}$ min$^{-1}$. In order that half life of both the processess is same, one should start with an initial concentration of `A' as `C' M. The value of `C' is.
\qn{4} Find the number of oxides of halogens exist only at low temperature.\\
$\mathrm{Br_2O}$, $\mathrm{ClO_2}$, $\mathrm{Cl_2O}$, $\mathrm{BrO_2}$, $\mathrm{BrO_3}$, $\mathrm{Cl_2O_7}$, $\mathrm{I_2O_7}$, $\mathrm{I_2O_5}$
\qn{5} Which of the following compounds form $2^\circ$--amines on reaction with $\mathrm{LiAlH_4}$
\begin{qoptions}
\opt{a}{\qfigi{p2_chem_s3_q05_1}}
\opt{b}{\qfigi{p2_chem_s3_q05_4}}
\opt{c}{$\mathrm{CH_3{-}CH_2{-}N{\equiv}C}$}
\opt{d}{$\mathrm{CH_3{-}CH_2{-}CN}$}
\opt{e}{\qfigi{p2_chem_s3_q05_3}}
\opt{f}{\qfigi{p2_chem_s3_q05_2}}
\opt{g}{\qfigi{p2_chem_s3_q05_5}}
\end{qoptions}
\qn{6} Two containers A and B of equal volume maintained at $400$ K and $600$ K are connected with a narrow tube of negligible volume as shown in figure.
\qfig{p2_chem_s3_q06_1}
Find number of moles of $\mathrm{H_2}$ gas in container A, when the stop cock is opened.
\partsec{3}{MATHEMATICS}
\sect{Section-I}
\qn{1} $\displaystyle\int\dfrac{\sec^2x}{\sqrt{a\sec^2x-b\tan^2x}}\,dx$ is -
\begin{qoptions}
\opt{A}{$\dfrac{1}{\sqrt{b-a}}\sin^{-1}\!\left(\tan x\sqrt{\dfrac{b-a}{a}}\right)+c$ if $b>a>0$}
\opt{B}{$\dfrac{1}{\sqrt{b-a}}\log_2\!\left(\tan x\sqrt{b-a}+\sqrt{a\sec^2x-b\tan^2x}\right)+c$ if $b>a>0$}
\opt{C}{$\dfrac{1}{\sqrt{a-b}}\sin^{-1}\!\left(\tan x\sqrt{\dfrac{a-b}{a}}\right)+c$ if $a>b>0$}
\opt{D}{$\dfrac{1}{\sqrt{a-b}}\log_e\!\left(\tan x\sqrt{a-b}+\sqrt{a\sec^2x-b\tan^2x}\right)+c$ if $a>b>0$}
\end{qoptions}
(where $c$ is integration constant)
\qn{2} If there exists a circle `$C$' whose one diameter is the join of focii of ellipse $E:\dfrac{x^2}{a^2}+\dfrac{y^2}{b^2}=1$ $(a>b)$ and other diameter is the minor axis of $E$, then -
\begin{qoptions}
\opt{A}{eccentricity of $E$ is $\dfrac12$}
\opt{B}{eccentricity of $E$ is $\dfrac{1}{\sqrt2}$}
\opt{C}{If $E$ is $\dfrac{x^2}{4}+\dfrac{y^2}{b^2}=1$ then director circle of circle `$C$' is $x^2+y^2=8$}
\opt{D}{If $E$ is $\dfrac{x^2}{4}+\dfrac{y^2}{b^2}=1$ then director circle of circle `$C$' is $x^2+y^2=4$}
\end{qoptions}
\qn{3} Let $p,q,r$ be distinct positive real numbers, such that exactly any two of the equations $px^2+2qx+r=0$; $qx^2+2rx+p=0$; $rx^2+2px+q=0$ have a common root, say $\alpha$, then which of the following statements is/are true -
\begin{qoptions}
\opt{A}{$\alpha$ is real \& negative}
\opt{B}{$\alpha$ is real \& positive}
\opt{C}{remaining third quadratic equation has real roots}
\opt{D}{remaining third quadratic equation has imaginary roots}
\end{qoptions}
\qn{4} The equation $a\cos3x+b\cos2x+c\cos x+\sin x=0$ $(a,b,c\in\mathbb R)$ has at least one root in $(0,2\pi)$, for
\begin{qoptions}
\opt{A}{$a=b=c=0$}
\opt{B}{$a+b+c=0$}
\opt{C}{$a$ is positive and $b$ \& $c$ must be of opposite sign}
\opt{D}{all values of $a,b,c$}
\end{qoptions}
\qn{5} Let `$O$' be the origin and the line $y=\left(2-\sqrt3\right)x$ cuts the line $x=1$ at point $A$. Right isosceles triangle $OAB$ is constructed such that hypotaneous is $OA$ and $B$ lies in first quadrant. Which of the following statements is (are) true ?
\begin{qoptions}
\opt{A}{length $AB=\sqrt3-1$}
\opt{B}{length $OB=\sqrt3+1$}
\opt{C}{$B\equiv\left(\dfrac{\sqrt3-1}{2},\dfrac{3-\sqrt3}{2}\right)$}
\opt{D}{$B\equiv\left(\dfrac{\sqrt3+1}{2},\dfrac{3+\sqrt3}{2}\right)$}
\end{qoptions}
\qn{6} If $A_1,A_2,\ldots\ldots A_n$ are $n$-distinct odd natural numbers not divisible by any prime greater than $5$, then $\dfrac1{A_1}+\dfrac1{A_2}+\dfrac1{A_3}+\ldots\ldots+\dfrac1{A_n}$ can never exceed -
\begin{qoptions}
\opt{A}{$1$}
\opt{B}{$2$}
\opt{C}{$3$}
\opt{D}{$4$}
\end{qoptions}
\sect{Section-II}
\qn{1} If $ax+\displaystyle\int_0^x\ln t\,dt=\dfrac{x^2}{3},\ \left(x\in\mathbb R^+\right)$ has exactly one solution, then the value of $(30(a+\ln3))$ is equal to
\qn{2} The tangent at point $P$ of a curve $y=f(x)$ meets the $y$-axis at $A$ and the line parallel to $y$-axis through $P$ meets the $x$-axis at $B$ such that area of triangle $OAB$ is $4$ square units ($O$ being the origin) and $f(1)=6$ \& $f(4)=9$. If distance between point of maxima and minima on $y=f(x)$ is $\sqrt d$, then the value of `$d$' is equal to
\qn{3} Let $f(x)=(x+1)(x+2)(x+3)(x+4)+5$ where $x\in[-6,6]$. If the range of the function is $[a,b]$ where $a,b\in\mathbb N$ then find the value of $(a+b)$.
\qn{4} If $P$ be point of intersection of line $\dfrac{x+1}{2}=\dfrac{y+2}{3}=\dfrac{z+3}{4}$ \& plane $x+y+4z=6$ \& $A,B,C$ are point of intersection of given plane with $X,Y$ and $Z$-axis respecitvely then find the distance of $P$ from centroid of $\triangle ABC$.
\qn{5} The polynomial $P(x)=(1+x+x^2+\ldots\ldots+x^{17})^2-x^{17}$ has $34$ complex roots of the form $z_k=r_k[\cos(2\pi a_k)+i\sin(2\pi a_k)]$, $k=1,2,3,\ldots..34$ with $0<a_1\le a_2\le a_3\le\ldots\ldots\le a_{34}<1$ and $r_k>0$. Given that $a_1+a_2+a_3+a_4+a_5=\dfrac pq$, where $p,q$ are relatively prime positive integers, then $(p+q)$ is equal to
\qn{6} A derivable function $f:\mathbb R^+\to\mathbb R$ satisfies the condition $f(x)-f(y)\ge\ln\dfrac xy+x-y\ \forall\ x,y\in\mathbb R^+$. If $g$ denotes the derivative of $f$ then compute the value of the sum $\displaystyle\sum_{n=1}^{100}g\!\left(\dfrac1n\right)$.
\sect{Section-III}
\qn{1} In a $\triangle ABC$ (with usual notations), if $a\cos B-b\cos A=\dfrac35c$, then $\dfrac{\tan A}{\tan B}$ is
\qn{2} The value of $k(k\in\mathbb R^+)$ so that area of the bounded region enclosed between the parabolas $y=x-kx^2$ and $y=\dfrac{x^2}{k}$ is maximum, is
\qn{3} In a $5$ team tournament, each team plays one game with every other team. The probability of winning any game by any team is $\dfrac12$ and there are no ties. Let $\dfrac mn$ be the probability that at the end of tournament there is no team with all wins or all defeats, $m$ and $n$ are relatively prime numbers. Then value $(2m-n)$ is
\qn{4} If the normal at $P(8,2)$ on the curve $xy=16$ meets the curve again at $Q$. The angle subtented by $PQ$ at origin is $\pi-\tan^{-1}\!\left(\dfrac{\lambda_1}{\lambda_2}\right)$ where $\lambda_1$ \& $\lambda_2$ are relatively prime, then $\lambda_1-2\lambda_2=$
\qn{5} Let $A$ be the matrix of order $2\times2$ and commute with $B=\begin{bmatrix}0&1\\0&0\end{bmatrix}$. Number of matrices $A$, if its elements belong to the set $S=\{-1,0,1\}$, is
\qn{6} If $f(x)=\displaystyle\int\dfrac{3(\ln x)^2+(\ln x)^4}{x\left(1+(\ln x)^2-(\ln x)^3\right)^2}\,dx$ and $f(1)=0$, then the value of $|3f(e^2)|$ is equal to
\end{multicols}
\end{document}
